\subsection{DEFINITION OF SEMI-SIMPLE LIE ALGEBRAS}

DEFINITION 1. Let \(\mathfrak{g}\) be a Lie algebra. \(\mathfrak{g}\) is called semi-simple if and only if the only commutative ideal of \(\mathfrak{g}\) is \(\{0\}\) .

Remarks. (1) The algebra \(\{0\}\) is semi-simple. An algebra of dimension 1 or 2 is not semi-simple (cf.~§ 5, no. 1, Example 1). There exist semi-simple algebras of dimension 3 (cf.~no. 7).

\begin{enumerate}
\def\labelenumi{(\arabic{enumi})}
\setcounter{enumi}{1}
\item
  A semi-simple algebra has zero centre and hence its adjoint representation is faithful.
\item
  If \(g_{1}, \ldots, g_{n}\) are semi-simple, \(g = g_{1} \times \cdots \times g_{n}\) is semi-simple: for if a is a commutative ideal of g, the projections of a onto \(g_{1}, \ldots, g_{n}\) reduce to \(\{0\}\) .
\end{enumerate}

THEOREM 1. Let \(\mathfrak{g}\) be a Lie algebra. The following conditions are equivalent:

\begin{enumerate}
\def\labelenumi{(\alph{enumi})}
\item
  g is semi-simple.
\item
  The radical r of g is zero.
\item
  The Killing form \(\beta\) of \(\mathfrak{g}\) is non-degenerate.
\end{enumerate}

Moreover, a semi-simple Lie algebra is equal to its derived ideal.

\begin{enumerate}
\def\labelenumi{(\alph{enumi})}
\item
  \(\Rightarrow\) (b): for if \(\mathfrak{r} \neq \{0\}\) , the last non-zero derived algebra of \(\mathfrak{r}\) is a commutative ideal of \(\mathfrak{g}\) .
\item
  \(\Rightarrow\) (c): this follows from Proposition 5 (b) of § 5, no. 5 (which at the same time proves the last assertion of the theorem).
\item
  \(\Rightarrow\) (a): this follows from Proposition 6 (b) of §4, no. 4.
\end{enumerate}

COROLLARY. Let \(\mathfrak{g}\) be a semi-simple Lie algebra and \(\rho\) a representation of \(\mathfrak{g}\) on a finite-dimensional vector space V. Then \(\rho(\mathfrak{g}) \subset \mathfrak{sl}(V)\) .

The linear form \(x \mapsto \operatorname{Tr} \rho(x)\) ( \(x \in \mathfrak{g}\) ) is zero when \(x\) is of the form \([y, z]\) ( \(y \in \mathfrak{g}, z \in \mathfrak{g}\) ) and hence on \(\mathcal{D}\mathfrak{g} = \mathfrak{g}\) .

PROPOSITION 1. Let \(\mathfrak{g}\) be a semi-simple Lie algebra and \(\rho\) a finite-dimensional faithful representation of \(\mathfrak{g}\) . Then the bilinear form on \(\mathfrak{g}\) associated with \(\rho\) is non-degenerate.

The orthogonal of g with respect to this form is a solvable ideal (§ 5, no. 4, Theorem 2) and hence is zero.

COROLLARY 1. Let g be a Lie algebra, \(\beta\) its Killing form and \(\mathfrak{a}\) a semi-simple subalgebra of g. The orthogonal \(\mathfrak{b}\) of \(\mathfrak{a}\) with respect to \(\beta\) is a supplementary subspace of \(\mathfrak{a}\) in g and \({[}\mathfrak{a}, \mathfrak{b}{]} \subset \mathfrak{b}\). If \(\mathfrak{a}\) is an ideal of g, so is \(\mathfrak{b}\), which is then the centralizer of \(\mathfrak{a}\) in g.

Let \(\beta'\) be the restriction of \(\beta\) to \(a\) : it is the bilinear form associated with the representation \(x \mapsto \operatorname{ad}_{\mathfrak{g}} x\) of \(a\) on the space \(\mathfrak{g}\) . This representation is faithful and hence \(\beta'\) is non-degenerate (Proposition 1). Hence \(\mathfrak{b}\) is supplementary to \(a\) in \(\mathfrak{g}\) . On the other hand, if \(x, y\) are in \(a\) and \(z \in \mathfrak{b}\) , then

\[
\beta (x, [ y, z ]) = \beta ([ x, y ], z) = 0,
\]

 since \([x,y] \in \mathfrak{a}\) , and hence \([y,z] \in \mathfrak{b}\) , which proves that \([\mathfrak{a},\mathfrak{b}] \subset \mathfrak{b}\) . If \(\mathfrak{a}\) is an ideal of \(\mathfrak{g}\) , we know that \(\mathfrak{b}\) is an ideal of \(\mathfrak{g}\) (§ 3, Proposition 7) and \(\mathfrak{g}\) is identified with \(\mathfrak{a} \times \mathfrak{b}\) . As the centre of \(\mathfrak{a}\) is zero, the centralizer of \(\mathfrak{a}\) in \(\mathfrak{g}\) is \(\mathfrak{b}\) .

COROLLARY 2. Every extension of a semi-simple Lie algebra by a semi-simple Lie algebra is semi-simple and trivial.

This follows immediately from Corollary 1.

COROLLARY 3. If g is semi-simple, every derivation of g is inner.

ad g is isomorphic to g and hence semi-simple and is an ideal of the Lie algebra d of derivations of g (§ 1, Proposition 1). If D ∈ d commutes with the elements of ad g, then, for all x ∈ g, ad D(x) = {[}D, ad x{]} = 0, whence D(x) = 0; hence D = 0. Corollary 3 then follows from Corollary 1.
% semantic-fragments: legacy-input-seam
\relax{}
